\documentclass[11pt]{article}
\usepackage[margin=1in]{geometry}
\usepackage{amsmath,amssymb,amsthm}
\usepackage[hidelinks]{hyperref}

\newtheorem{lemma}{Lemma}
\newtheorem{proposition}{Proposition}
\theoremstyle{definition}
\newtheorem{definition}{Definition}
\theoremstyle{remark}
\newtheorem{remark}{Remark}

\title{Kobon triangles beyond simple arrangements:\\ a complete SAT model and the case of eleven lines\\[4pt]
\large Working draft, version 0 (27 October 2026)}
\author{Papex Labs}
\date{}

\begin{document}
\maketitle

\begin{abstract}
The Kobon triangle problem asks for the largest number $K(n)$ of bounded triangular faces in an
arrangement of $n$ lines in the plane. The classical bound $K(n) \le \lfloor n(n-2)/3 \rfloor$ is
proved by counting segments, an argument that assumes every segment is a side of at most one
triangle. That assumption holds in simple arrangements but fails when three or more lines meet at a
point. In October 2026 the values $a(11)$ to $a(17)$ were removed from OEIS A006066 after a
counterexample to the lemma that extended the argument to all arrangements. We describe a SAT model
whose solutions are exactly the affine arrangements of $n$ pseudolines, simple or not, together with
a set of their bounded triangular faces, and we validate it by reproducing $K(n)$ for $3 \le n \le 7$
over all arrangements. The model is designed to decide whether some arrangement of $11$ lines,
simple or not, has $33$ bounded triangles. This version reports the model, its validation and the
status of the computation; it states no result for $n = 11$ yet.
\end{abstract}

\section{The question}

For $n$ lines in the real affine plane, let $K(n)$ be the largest number of bounded faces that are
triangles. Honma found an arrangement of $11$ lines with $32$ triangles, and the counting bound gives
$33$. For simple arrangements of $11$ pseudolines meeting the bound (\emph{perfect} arrangements),
Savchuk~\cite{savchuk} showed by SAT that $33$ triangles are impossible. The reduction of the general
case to simple perfect arrangements rested on Lemmas~1 and~2 of Cl\'ement and Bader~\cite{cb}. In
2026 Parpalak and Utkin pointed out a counterexample to Step~3 of Lemma~1: in the arrangement formed
by the sides and medians of a triangle, the centroid is a vertex of six pairs of triangles sharing a
side~\cite{oeis}. On 9 October 2026 Savchuk removed $a(11)$ to $a(17)$ from OEIS A006066 and recorded
this reason~\cite{oeis}.

The draft of Cl\'ement and Bader states that Tamura's bound ``directly follows from the proof of
Lemma~1''~\cite{cb}. So for arrangements in which three or more lines meet at a point, both the
value $a(11)$ and the bound $33$ are open as far as these sources go.

\section{Why non-simple arrangements escape the counting argument}

A \emph{vertex} is a point where at least two lines meet; its multiplicity $m_v$ is the number of
lines through it. A \emph{segment} is a bounded piece of a line between two consecutive vertices
on it. Each side of a bounded triangular face is exactly one segment: a line through an interior
point of a side would cross into the triangle.

\begin{lemma}\label{lem:simple}
If a segment $PQ$ is a side of two triangular faces, then $m_P \ge 3$ or $m_Q \ge 3$.
\end{lemma}
\begin{proof}
Let $PQ$ lie on line $L$. A triangle on one side of $L$ is bounded by $L$, a line $A$ through $P$
and a line $B$ through $Q$, with apex $A \cap B$ on that side; a triangle on the other side uses
lines $A'$ through $P$ and $B'$ through $Q$, with apex on the other side. If $A = A'$ and $B = B'$
then $A$ and $B$ meet on both sides of $L$, which two lines (or pseudolines) cannot do. Hence
$A \ne A'$ or $B \ne B'$, and $P$ or $Q$ lies on at least three lines.
\end{proof}

In a simple arrangement each segment is therefore a side of at most one triangle, there are at most
$n(n-2)$ segments, and $3T \le n(n-2)$. Without simplicity, write $S$ for the number of segments and
$D$ for the number of segments that are sides of two triangles. Then $3T \le S + D$. A vertex of
multiplicity $m$ removes $m(m-2)$ segments compared with a simple vertex set, a class of $q$ parallel
lines removes $q(q-1)$, and by Lemma~\ref{lem:simple} each doubly used segment can be charged to an
endpoint of multiplicity at least $3$, which is an endpoint of at most $2m$ segments. This gives
\[
  3T \;\le\; n(n-2) \;-\; \sum_{\text{parallel classes}} q(q-1) \;-\; \sum_{v} m_v(m_v-4),
\]
where the last sum runs over vertices with $m_v \ge 3$. A vertex of multiplicity $3$ contributes
$+3$, so this argument alone does not even give $3T \le n(n-2)$. The triangle with its medians
shows that doubly used segments do occur.

\section{A model over all pseudoline arrangements}

\paragraph{Arrangements as oriented matroids.} An affine arrangement of $n$ pseudolines is a
projective arrangement of $n+1$ pseudolines with one of them, labelled $0$, declared to be the line
at infinity; parallel lines are lines meeting on line $0$. By the Folkman--Lawrence topological
representation theorem~\cite{fl,bjorner}, projective arrangements of pseudolines correspond to simple
rank-$3$ oriented matroids. We encode a chirotope $\chi$ on the triples of $\{0,\dots,n\}$ with values
in $\{+,-,0\}$ and impose:
\begin{itemize}
\item $\chi$ is alternating;
\item every pair of elements lies in a triple with $\chi \ne 0$ (distinct pseudolines);
\item the zero triples form a simple rank-$3$ matroid: if $abc$ and $abd$ are dependent, so are
  $acd$ and $bcd$;
\item the $3$-term Grassmann--Pl\"ucker relations: for every $a$ and $x_1,\dots,x_4$, the products
  $\chi(a,x_1,x_2)\chi(a,x_3,x_4)$, $-\chi(a,x_1,x_3)\chi(a,x_2,x_4)$, $\chi(a,x_1,x_4)\chi(a,x_2,x_3)$
  are all zero or include both signs.
\end{itemize}
A support that is a matroid together with the $3$-term relations characterises chirotopes
\cite[Thm.~3.6.2]{bjorner}.

\paragraph{Bounded triangles.} For non-parallel lines $a,b$, the vertex $a \cap b$ lies on side
$s_{ab}(d) = \chi(a,b,d)\,\chi(a,b,0)$ of a line $d$; the factor $\chi(a,b,0)$ normalises the vertex by
the line at infinity. Lines $a,b,c$ bound a triangular face exactly when they are pairwise
non-parallel, $\chi(a,b,c) \ne 0$, and no other line $d$ gives two of the three vertices opposite
nonzero sides: such a $d$ would meet the open triangle, and a $d$ that meets it must do so.

\paragraph{Encoding.} Variables encode $\chi$, the vertex sides and an indicator $T_{abc}$ for each
triple of lines; $T_{abc}$ implies the face conditions, and a sequential counter requires at least
$K$ indicators to be true. Only this direction is needed: an arrangement with $K$ triangles gives a
satisfying assignment, so an unsatisfiable instance excludes all arrangements of $n$ pseudolines,
hence of $n$ lines, with $K$ bounded triangles. Symmetry breaking fixes $T_{123}$ (relabelling) and
the orientation of each element ($\chi(0,1,2)=+$, $\chi(0,1,e)\in\{+,0\}$ for $e\ge 3$,
$\chi(0,2,e)=+$ when $\chi(0,1,e)=0$, and $\chi(1,2,3)=+$); every reorientation class contains an
element meeting these conditions.

\section{Validation}

For $3 \le n \le 7$ the model is satisfiable at $K = K(n)$ and unsatisfiable at $K = K(n)+1$, with
$K(n) = 1, 2, 5, 7, 11$, matching OEIS A006066. These runs cover non-simple arrangements, so they
confirm the values for all arrangements of up to $7$ pseudolines. Runs for $n = 8, 9, 10$ and the
simple case of $n = 11$ (which must reproduce Savchuk's result) are in progress.

\section{Status and plan}

The instance for $n = 11$, $K = 33$ over all arrangements is the target. If it is unsatisfiable,
an LRAT or DRAT certificate checked by an independent checker will establish $a(11) = 32$ and the
bound for $n = 11$. If it is satisfiable, the solution is a pseudoline arrangement with $33$ bounded
triangles, which must then be straightened, or shown non-stretchable, and checked exactly. Sound
redundant constraints from Section~2 (no vertex of multiplicity above $4$ unless compensated,
charging of doubly used segments) may be needed for the solver.

\section*{How this work was done}
This draft, the model and the code were prepared by Claude Opus 5.5, an AI model, working for Papex
Labs; a person at Papex Labs reviews and signs each version. The code was written from the published
statements of the sources, not from the tools of~\cite{savchuk}. Solvers: Kissat 4.0.4 and CaDiCaL
3.0.1. Code and results will be released with the first version that states a result.

\begin{thebibliography}{9}
\bibitem{bjorner} A.~Bj\"orner, M.~Las Vergnas, B.~Sturmfels, N.~White, G.~M.~Ziegler,
  \emph{Oriented Matroids}, 2nd ed., Cambridge University Press, 1999.
\bibitem{cb} G.~Cl\'ement and J.~Bader, \emph{Tighter upper bound for the number of Kobon
  triangles}, unpublished draft, 2007. \url{https://oeis.org/A006066/a006066.pdf}
\bibitem{fl} J.~Folkman and J.~Lawrence, Oriented matroids, \emph{J. Combin. Theory Ser. B} 25
  (1978) 199--236.
\bibitem{oeis} OEIS Foundation, Sequence A006066, \url{https://oeis.org/A006066}, comments and
  edit history as of 27 October 2026.
\bibitem{savchuk} P.~Savchuk, Constructing optimal Kobon triangle arrangements via table encoding,
  SAT solving, and heuristic straightening, arXiv:2507.07951v2, 2026.
\end{thebibliography}
\end{document}
